\section{Motivation for a Structured Approach}

Part 1's architectural problems shared a common origin, and the conclusion drawn from them shaped everything that followed. The missing service layer, the \textbf{\texttt{MainPage.js}} monolith, the scattered reads of authentication state, the unauthenticated \textbf{\texttt{user\_id}} parameter, and the test that was made to pass without its underlying inconsistency being addressed were not, individually, failures of the agent's capability. Each decision was reasonable in isolation; what was missing in every case was something persistent to check the decision against --- an existing convention to extend, a plan to follow, a review step that had to happen before the code was accepted rather than only once a problem had already spread. The conclusion drawn from this was explicit: the quality of the software produced in Part 1 had been shaped less by the coding model itself than by the architectural constraints, review process, and development rules imposed on the agent --- or, in Part 1's case, by their absence.

This reframed the question for the second part of the project. Rather than treating better output as something to get by giving the agent more capability or more detailed individual prompts, Part 2 was designed around the idea that the missing ingredient was structure: a defined set of rules the agent had to follow regardless of the feature it was working on, a point at which a plan had to be reviewed before any code was written, a role dedicated to checking the result afterward rather than relying on a problem being noticed by chance, and an explicit requirement to determine why a test had failed before changing anything, rather than defaulting to whichever fix made the test pass.

Realizing this structure required an environment built to support it, which motivated moving development from Claude Code inside Visual Studio Code to Google Antigravity --- an agent-first environment organized around coordinating multiple agents rather than a single conversational thread, with native support for agent management, planning, parallel agents, and project context that persists across sessions rather than needing to be re-supplied in every prompt. What this made possible --- an AI Development Charter defining architectural rules before implementation began, a division of labor between an agent responsible for planning and one responsible for implementing, a mandatory design review gate before coding, a dedicated review role after it, and a formalized testing discipline --- is the subject of the sections that follow.

